%%
%% Copyright 2019-2020 Elsevier Ltd
%%
%% This file is part of the 'CAS Bundle'.
%% --------------------------------------
%%
%% It may be distributed under the terms of the LaTeX Project Public
%% License, either version 1.2 of this license or (at your option) any
%% later version.  The latest version of this license is in
%%    http://www.latex-project.org/lppl.txt
%% and version 1.2 or later is part of all distributions of LaTeX
%% version 1999/12/01 or later.
%%
%% The list of all files belonging to the 'CAS Bundle' is
%% given in the file `manifest.txt'.
%%
%% Template article for cas-sc documentclass for
%% double column output.

%\documentclass[a4paper,fleqn,longmktitle]{cas-sc}
\documentclass[a4paper,fleqn]{cas-sc}
\usepackage{njfutitlepage}
\usepackage[authoryear]{natbib}
\usepackage{csquotes}
\usepackage[UTF8]{ctex}
\usepackage[ruled,vlined]{algorithm2e}
\usepackage{etoolbox}
\usepackage{placeins}
\usepackage{booktabs}
\usepackage{gbt7714}

\newcommand{\refSec}[1]{\hyperref[#1]{Section~\ref*{#1}}}
\newcommand{\refFig}[1]{\hyperref[#1]{Fig.~\ref*{#1}}}
\newcommand{\refTab}[1]{\hyperref[#1]{Table~\ref*{#1}}}
\renewcommand{\figurename}{Fig.}

\newcommand{\FeHALSRepo}{\href{https://github.com/WZYivan/FeHALS}{FeHALS github Main Repo}}

%%%Author definitions
\def\tsc#1{\csdef{#1}{\textsc{\lowercase{#1}}\xspace}}
\tsc{WGM}
\tsc{QE}
\tsc{EP}
\tsc{PMS}
\tsc{BEC}
\tsc{DE}
%%%

\begin{document}
\let\WriteBookmarks\relax
\def\floatpagepagefraction{1}
\def\textpagefraction{.001}

\shortauthors{7组}
\shorttitle{}
\title[mode=title]{\makecell[c]{FeHALS：基于Web的3D可视化航路规划与激光仿真系统\\\textbf{技术设计书}}}

\Course{\enquote{智能+}GIS设计与开发}
\Major{测绘工程}
\Grade{2023级}
\Group{7 组}

\Member[2310604116]{王 孜悠}
\Member[2310604109]{梅 康凯}
\Member[2310604115]{唐 如意}
\Member[2310604122]{杨 杰瑞}
\Member[2310604117]{温 璞}
\Member[2110604112]{刘 继明}
\Member[2110604106]{李 英杰}
\Member[2310604103]{杜 仲宇}
\Member[2310604113]{施 祺}
\Member[2310604118]{向 宇鸣}

\Teacher{隋 铭明}
\Teacher{那 嘉明}
\Teacher{陈 动}
\Teacher{朱 杰}

\MakeNJFUTitlePage[UseDocTitle]
\tableofcontents
\newpage

\section{背景目的}

\subsection{项目背景}

机载激光扫描（Airborne Laser Scanning，ALS）能够直接获取地表及地物的三维离散采样，是数字高程模型生产、森林调查、城市三维建模和灾害监测的重要数据源。真实飞行采集受设备成本、空域、天气和任务风险制约，外业前通过虚拟场景试验航线与扫描参数，可以降低反复试飞的成本。

HELIOS++是目前最为活跃的开源激光雷达仿真框架之一，采用CPU光线追踪模拟激光脉冲从发射到返回的全波形响应，支持地基、车载、无人机和机载平台的扫描配置\citep{heliosPlusPlus,heliosPlusPlus2}。其核心架构将场景（scene）、平台（platform）、扫描器（scanner）和测量任务（survey）四个维度解耦，以XML文件定义各组件间的组合关系。平台支持线性路径（linear path）和插值运动（interpolated）两种运动模型，其中插值运动模型通过读取时间戳序列的轨迹文件（.trj），在相邻航点间线性插值获得平台位置与姿态，再根据ARINC 705规范将偏航角（yaw/heading）转换为绕Z轴的旋转，以驱动扫描器产生与航线正交的扫描线。这一设计使HELIOS++能够在不依赖真实飞控数据的情况下，仅凭航点坐标和偏航角序列即可生成物理上合理的激光扫描模拟。

然而，HELIOS++的原生操作方式为命令行和XML配置，用户需手动编写场景、平台、扫描器、survey等多个XML文件，并管理轨迹文件与资产的路径引用。对于初学者或测绘工程专业的学生而言，这一过程存在较高的入门门槛：参数的单位与范围不直观、轨迹与配置的关联关系隐含在XML索引中、仿真结果缺乏即时的可视化反馈。此外，HELIOS++官方提供的QGIS插件AEOS\citep{aeos}虽在一定程度上降低了操作门槛，但AEOS依赖于QGIS桌面环境，安装配置复杂，且在使用过程中仍需用户具备一定的QGIS操作经验。

FeHALS（Frontend of HELIOS++ for Airborne Laser Scanning）正是在这一背景下提出的，它以浏览器为操作入口，将三维预览、航点编辑、弓字形航线生成、参数配置、HELIOS++任务执行、实时日志和点云分析串联为完整流程，旨在为HELIOS++提供一个现代化、免配置的Web前端套件。开放式WebGIS技术便于跨平台访问与后续扩展\citep{opengis,viswebdrone,lidarWebgis}，浏览器侧利用WebGL和Three.js完成硬件加速三维显示，无需安装专用桌面GIS软件。在GIS应用层面，FeHALS将激光雷达仿真中的坐标系、高程模型、航线规划、点云统计分析等GIS核心概念融入Web工作流，使仿真过程与GIS工作流自然衔接\citep{pointCloudWeb}。

\subsection{建设目标}

本项目旨在为HELIOS++提供直观、可交互、可追踪的前端工作台：集中管理模型、航迹和扫描参数；在Z-up三维场景中完成航点规划；生成HELIOS++所需轨迹和XML；以异步任务运行仿真并实时显示日志；解析XYZ、LAS、LAZ结果，提供点云渲染与统计；形成便于多人协作和迭代的WebGIS工程结构。

系统面向激光雷达仿真学习者、测绘工程学生和方案验证人员，定位为教学与原型验证工具，不替代实际航空作业中的空域审批、飞控、坐标转换及成果质量验收。本设计书以GitHub仓库master分支当前实现为准，规划中的LOD、任务队列、Docker等功能不会当作现成功能描述。

\subsection{设计原则}

系统遵循坐标与参数一致、操作可视化、任务可追踪、前后端模块化和能力渐进增强原则。前端、后端和HELIOS++均以Z轴表示高程；任务具有唯一标识、状态、日志和结果；组件、状态、接口和服务分层；当前先满足课程演示和百万点级预览，超大点云与生产部署作为后续建设内容。

\section{需求分析}

\subsection{业务流程}

一次完整业务流程为：启动前后端并检查HELIOS++环境；上传场景模型并检查方向和包围盒；手工编辑航点或由矩形生成弓字形航线；设置平台、飞行与扫描参数；生成轨迹和配置；提交仿真并查看实时日志；加载结果点云、查看统计与直方图；下载结果或清理缓存。

在仿真数据源方面，系统可处理两类典型的三维场景数据：一类是建筑群级别的模型集合，如BuildingWorld数据集中的单体建筑模型，这类模型细节丰富、结构复杂，适合用于验证仿真参数对不同建筑形态的响应差异；另一类是建筑内部结构模型，如Building3D数据集中的精细室内场景，其空间尺度小、遮挡关系复杂，适合用于评估激光穿透率和点云完整性。两类数据集对于验证FeHALS的航高自动计算、扫描角度配置和点云覆盖度分析等功能具有重要参考价值。

\subsection{功能需求}

系统功能需求如\refTab{tab:functions}所示。

\begin{table}[pos=!htbp]
    \centering
    \caption{FeHALS功能需求。}
    \begin{tabular*}{\textwidth}{@{\extracolsep{\fill}} l l}
        \toprule
        编号 & 需求说明 \\
        \midrule
        FR-01 & 上传OBJ、GLTF、GLB、STL；列出、显示/隐藏、删除模型，显示包围盒。仅OBJ作为当前HELIOS++场景输入。 \\
        FR-02 & 提供旋转、缩放、平移、视角适配、Z-up网格和模型方向转换。 \\
        FR-03 & 支持点击添加、拖动、右键或Delete删除、列表编辑和清空，并以折线连接。 \\
        FR-04 & 选择矩形范围和航线间距，自动生成连续往返的弓字形航点。 \\
        FR-05 & 以所有模型世界包围盒的最大Z加20m计算建议高度，并受平台范围约束。 \\
        FR-06 & 设置UAV/Airborne、速度、高度、扫描频率、扫描半角、脉冲频率和输出格式。 \\
        FR-07 & 生成恒定高度轨迹文件和HELIOS++ survey XML。 \\
        FR-08 & 异步调用HELIOS++，记录pending/running/completed/failed状态并支持取消。 \\
        FR-09 & WebSocket实时推送日志，REST接口可补取历史日志。 \\
        FR-10 & 查询、下载和解析结果；显示点数、边界、高程/强度统计及40组直方图。 \\
        FR-11 & 统计并分类清理models、trajectories、configs、results。 \\
        \bottomrule
    \end{tabular*}
    \label{tab:functions}
\end{table}

\subsection{与现有工具的对比分析}

FeHALS不是对现有点云编辑软件或仿真引擎的简单替代，而是定位于它们之间的桥梁：它既不是专业的点云编辑工具如CloudCompare\citep{cloudcompare}，也不是点云处理库如PDAL\citep{pdal}，更不是激光雷达仿真引擎HELIOS++本身，而是将HELIOS++的能力封装为Web可访问的图形化工作流。\refTab{tab:comparison}展示了FeHALS与同类工具的功能定位差异。

\begin{table}[pos=!htbp]
    \centering
    \caption{FeHALS与同类工具的功能定位对比。}
    \begin{tabular*}{\textwidth}{@{\extracolsep{\fill}} l l l l l}
        \toprule
        维度 & FeHALS & AEOS (QGIS Plugin) & CloudCompare & PDAL \\
        \midrule
        定位 & HELIOS++ Web前端 & HELIOS++ QGIS插件 & 点云编辑与可视化 & 点云处理库 \\
        安装方式 & 浏览器访问 & QGIS插件安装 & 桌面应用安装 & Python/C++库 \\
        三维场景编辑 & 交互式3D场景 & 无独立3D场景 & 专业3D编辑 & 无 \\
        航迹规划 & 航点+弓字形生成 & 需手动绘制 & 不支持 & 不支持 \\
        仿真驱动 & 一键执行 & 需切换至HELIOS++ & 不支持 & 不支持 \\
        点云分析 & 全量统计+直方图 & 有限 & 丰富 & 丰富 \\
        \bottomrule
    \end{tabular*}
    \label{tab:comparison}
\end{table}

FeHALS的核心竞争壁垒在于：将HELIOS++仿真引擎的复杂XML配置和命令行操作封装为直观的Web界面，同时保留对仿真参数的精细控制权。与AEOS相比，FeHALS不依赖QGIS桌面环境，安装配置成本更低；与CloudCompare和PDAL相比，FeHALS专注的是仿真前的规划和仿真后的即时分析，而非点云的专业编辑和处理。三者之间是互补关系——用户可以在FeHALS中完成仿真规划与初步分析，再将结果导出至CloudCompare或PDAL进行深度处理。

\subsection{非功能需求与约束}

\newcommand{\nsitem}[1]{\item[\textbf{#1}]}

\begin{description}
    \nsitem{易用性} 核心流程在单页完成，参数显示单位、范围和默认值；三维场景支持鼠标拖拽、滚轮缩放等直观操作；参数校验在提交前给出明确提示。
    \nsitem{性能} 全量统计与显示数据分离，前端渲染最多一百万点，XYZ每五十万行分块读取；统计计算基于全量点集，确保准确性。
    \nsitem{可靠性} 子进程失败不拖垮API，模型异步加载期间删除不产生残留对象；WebSocket断线后自动重连，日志可通过REST补取。
    \nsitem{可维护性} 视图、Pinia状态、API客户端、路由、数据模型和业务服务解耦；后端按API路由、Pydantic schema、业务service和config分层。
\end{description}

系统依赖正确安装的HELIOS++及其环境路径，未配置时只能使用三维预览。UAV与Airborne当前均引用copter\_linearpath平台和riegl\_vux-1uav扫描器，主要由参数范围区分。航点在$z=0$地面编辑，生成轨迹时统一叠加飞行高度，属于恒高简化模型。仓库说明记录本机构建的LASlib输出异常，因此默认XYZ；LAS/LAZ只有在HELIOS++正确链接LASlib时可用。任务和配置主要存于单进程内存，系统也未设置账户权限，生产部署前必须补充持久化、鉴权、上传限制、HTTPS和审计。

\section{总体设计}

\subsection{系统架构}

FeHALS采用B/S三层架构。与传统C/S架构的GIS系统相比，B/S架构具有免安装、跨平台、易升级等显著优势——用户无需安装任何桌面软件，仅通过浏览器即可访问整个系统，服务端更新后所有用户立即获得最新功能。这种架构在WebGIS领域已得到广泛验证和应用\citep{viswebdrone,lidarWebgis,pointCloudWeb}。

表现层基于Vue 3组合式API构建单页应用（SPA），利用Vite实现开发热更新和构建打包\citep{vite}。Vue 3的Composition API允许将逻辑按功能组织为可复用的组合式函数（Composables），例如useThreeScene封装所有Three.js场景逻辑，useHeliosAPI封装所有后端API调用。Pinia作为状态管理库\citep{pinia}，将导航状态（waypoints store）、场景状态（scene store）和仿真状态（simulation store）三者分离，避免了组件间跨层级的事件传递。Axios统一管理HTTP请求，支持请求拦截和响应转换。Three.js负责三维场景的渲染和交互，使用WebGL实现硬件加速\citep{threejs,webgl}。

服务层基于FastAPI框架\citep{fastapi}，利用Python asyncio实现异步处理。Pydantic对所有请求和响应进行类型校验，在业务逻辑执行前捕获格式错误\citep{pydantic}。业务逻辑按服务模块组织：模型管理（model\_service）、轨迹生成（trajectory\_generator）、配置生成（config\_generator）、仿真执行（simulation\_runner）和结果解析（result\_parser）。WebSocket端点在任务执行期间推送实时日志，REST接口用于查询任务状态和结果。

计算层即HELIOS++仿真引擎，通过asyncio子进程调用。系统不直接链接HELIOS++的C++库，而是通过命令行参数传递配置，通过标准输出捕获日志，通过文件系统获取结果，这种松耦合方式便于故障隔离和日志记录。

本地目录保存模型、轨迹、配置和结果，模型manifest用于服务重启后的恢复。前端只交换JSON、文件和任务标识，后端负责文件组织、配置生成、进程管理和结果解析。

系统总体架构如\refFig{fig:arch}所示。

\begin{figure}[pos=!htbp]
    \centering
    \includegraphics[width=0.85\textwidth]{figs/gernal_fehals_arch.pdf}
    \caption{FeHALS总体架构}
    \label{fig:arch}
\end{figure}

前端中Scene3D.vue管理画布，ControlPanel.vue管理参数，WaypointList.vue管理航点，LogConsole.vue展示日志，PointCloudPanel.vue展示结果。共享状态由scene、waypoints、simulation三个store管理；跨组件逻辑放在composables。后端按API路由、Pydantic schema、业务service和config分层。

系统模块架构如\refFig{fig:module}所示。

\begin{figure}[pos=!htbp]
    \centering
    \includegraphics[width=\textwidth]{figs/detialed_fehals_arch.pdf}
    \caption{FeHALS模块架构}
    \label{fig:module}
\end{figure}

\subsection{数据与控制流程}

组件操作先更新Pinia，再由useHeliosAPI请求后端。后端为轨迹、配置和任务分配UUID，生成文件并启动HELIOS++子进程；日志写入任务缓冲并广播；成功后定位结果文件、计算全量统计并返回显示抽样；前端创建THREE.Points绘制点云。

数据流程遵循清晰的单向依赖链：场景模型和航点数据是输入起点，轨迹文件和参数配置由输入数据生成，配置XML引用轨迹文件和场景文件，HELIOS++读取配置后执行仿真，仿真结果经解析后生成点云统计和渲染数据。这一流程确保了数据变更可以追踪，便于定位问题和迭代优化。

系统工作流程如\refFig{fig:workflow}所示。

\begin{figure}[pos=!htbp]
    \centering
    \includegraphics[width=0.85\textwidth]{figs/gernal_fehals_workflow.pdf}
    \caption{FeHALS工作流程}
    \label{fig:workflow}
\end{figure}

开发环境中Vite默认使用5173端口，并把/api和/ws转发到8000端口FastAPI；后端通过/static发布文件。Windows用run.bat，类Unix用run.sh。当前是单机双进程课程演示结构，尚未包括数据库、反向代理、容器编排和集中监控。

\subsection{关键设计决策}

采用Web前端可以降低分发门槛并把场景、参数与日志放在同一工作区\citep{viswebdrone,lidarWebgis}。在技术选型上，Vue 3被选为前端框架的核心原因在于其组合式API天然适合将三维场景、航点编辑、仿真控制等复杂逻辑拆分为独立的组合式函数，每个函数封装特定领域的逻辑和状态，互不干扰。Pinia相比Vuex具有更完善的TypeScript类型推断和更简洁的API，在需要深度追踪仿真任务状态变更时优势明显。Three.js作为WebGL的封装库，提供了一整套三维场景管理工具，无需从零编写着色器即可实现模型加载、航点标记、点云渲染等核心功能。

系统显式将相机up设为$(0,0,1)$，Y-up模型加载时绕X轴旋转$90^\circ$，保证测绘高程语义一致。资源操作采用REST，连续日志采用WebSocket\citep{fielding2000rest,WebSocket}。GLTF/GLB/STL只承担当下前端预览，避免与HELIOS++输入能力混淆。点云采取\enquote{全量统计、抽样渲染}，在统计准确性与交互性能间折中；更大数据仍需要层次LOD或渐进渲染\citep{pointCloudRendering,pointCloudTerrain}。

\section{功能设计}

\subsection{模型、场景与航点}

后端校验扩展名并以随机标识保存上传文件，返回静态URL。前端按扩展名选择OBJLoader、GLTFLoader或STLLoader，将对象加入modelsGroup。模型可见性和包围盒可切换，Box3计算尺寸。若模型在异步加载完成前已被删除，pendingRemoval会使迟到结果被释放，避免\enquote{幽灵对象}。相机按包围盒自适应，地面为XY平面，Z为高程。

Raycaster与$z=0$隐形地面求交：点击添加航点，拖动修改XY，右键或Delete删除，红色球体和折线表示路径，选中航点变黄。拖动时临时禁用相机控制，以25像素平方阈值区分点击和拖动。弓字形航线在矩形内按间距$d$生成，条数为

\begin{equation}
    N = \left| \frac{y_{\max} - y_{\min}}{d} \right| + 1,
\end{equation}

相邻扫描线的X端点次序交替，形成连续往返路径。

\subsection{参数、轨迹和配置}

建议高度取场景包围盒并集最高点$z_{\max}$加$20\text{ m}$：

\begin{equation}
    h_{\text{suggest}} = z_{\max} + 20\text{ m}.
\end{equation}

UAV速度为$0.5$--$50\text{ m/s}$，高度为$3$--$500\text{ m}$；Airborne速度为$10$--$200\text{ m/s}$，高度为$3$--$5000\text{ m}$。RIEGL VUX-1UAV扫描频率为$10$--$200\text{ Hz}$，扫描半角为$1^\circ$--$165^\circ$，脉冲频率为$50$--$550\text{ kHz}$。

轨迹文件（.trj）是HELIOS++插值运动模型的核心输入，其格式为每行包含七个以逗号分隔的字段：时间（t）、滚转角（roll）、俯仰角（pitch）、偏航角（yaw）、X坐标、Y坐标、Z坐标。系统在生成轨迹时，偏航角由相邻航点间的方向向量实时计算：

\begin{equation}
    \text{yaw} = \arctan\left(\frac{x_{i+1} - x_i}{y_{i+1} - y_i}\right).
\end{equation}

该公式符合ARINC 705规范，yaw=0°表示正北方向（+Y），yaw=90°表示正东方向（+X），顺时针方向增加。在航点转角处，系统通过生成同一时间戳、同一位置、不同偏航角的两行记录实现瞬时转向，使每个航段保持恒定偏航角，确保扫描线始终正交于航线方向。飞行速度通过目标速度（默认10 m/s）和航段距离计算时间步长，各航段时间连续累积。

后端将航点改写为统一高度的轨迹文件，并把平台、扫描器、轨迹、速度、频率、半角和格式写入survey XML。轨迹生成与配置生成分离，便于复用和排错。扫描条带宽度由飞行高度和扫描角决定：

\begin{equation}
    W = 2 \cdot h \cdot \tan\left(\frac{\theta}{2}\right),
\end{equation}

其中$h$为飞行高度，$\theta$为总扫描角。扫描线间距由飞行速度和扫描频率决定：

\begin{equation}
    S = \frac{v}{f_{\text{scan}}},
\end{equation}

其中$v$为飞行速度，$f_{\text{scan}}$为扫描频率。合理的参数配置应使扫描线间距小于或等于激光脚点间距，以保证点云覆盖的完整性。

\subsection{仿真、日志和结果}

任务提交后检查trajectory、config和可选scene model，分配UUID并由pending进入running。HELIOS++标准输出和错误输出逐行保存并广播到/ws/logs/\{task\_id\}；退出码为0且找到结果时进入completed，否则进入failed。运行中任务可取消。前端保存taskId、状态、进度、日志和结果，WebSocket不可用时可通过REST查询。

系统还提供HELIOS++运行环境诊断功能，在设置面板中集中展示检测状态。诊断项包括：HELIOS++可执行文件的存在性和可执行权限、HELIOS++资源目录的完整性（平台定义、扫描器定义、场景资产搜索路径）、静态工作目录的就绪状态。诊断结果以状态徽章和明细列表的形式展示，缺失项给出引导提示，帮助用户快速定位环境配置问题。

LAS/LAZ由laspy读取\citep{laspy}，XYZ分块读取并过滤非法记录。统计包括总点数、XYZ边界、高程均值、标准差、最小/最大、中位数、P5/P95、40组高程直方图，以及存在强度字段时的强度统计。渲染超过一百万点时使用覆盖全文件序列的等间隔索引抽样。前端以BufferGeometry和PointsMaterial绘制，支持固定色、高度色和强度色；归一化为

\begin{equation}
    t = \frac{v - v_{\min}}{\max(v_{\max} - v_{\min}, 1)},
\end{equation}

再映射至蓝—青—绿—黄—红色带。抽样仅影响显示，不改变全量统计。

\subsection{接口设计}

系统后端REST接口如\refTab{tab:apis}所示。

\begin{table}[pos=!htbp]
    \centering
    \caption{FeHALS后端接口。}
    \begin{tabular*}{\textwidth}{@{\extracolsep{\fill}} l l l}
        \toprule
        方法 & 路径 & 说明 \\
        \midrule
        POST & /api/models/upload & 上传模型 \\
        GET & /api/models & 查询模型列表 \\
        DELETE & /api/models/\{model\_id\} & 删除模型 \\
        POST & /api/trajectory/generate & 生成航迹 \\
        GET & /api/trajectories & 查询航迹列表 \\
        GET & /api/trajectories/download/\{file\_id\} & 下载航迹 \\
        POST & /api/config/generate & 生成配置 \\
        POST & /api/simulation/run & 创建仿真任务 \\
        GET & /api/simulation/status/\{task\_id\} & 查询任务状态 \\
        GET & /api/simulation/logs/\{task\_id\} & 获取任务日志 \\
        POST & /api/simulation/cancel/\{task\_id\} & 取消任务 \\
        GET & /api/results/\{task\_id\} & 获取结果统计和抽样 \\
        GET & /api/results/\{task\_id\}/download & 下载结果 \\
        GET & /api/cache & 查看缓存统计 \\
        DELETE & /api/cache/\{cache\_type\} & 分类清理缓存 \\
        GET & /api/env/diagnose & 环境诊断 \\
        \bottomrule
    \end{tabular*}
    \label{tab:apis}
\end{table}

主要数据对象如\refTab{tab:dataobjects}所示。

\begin{table}[pos=!htbp]
    \centering
    \caption{主要数据对象。}
    \begin{tabular*}{\textwidth}{@{\extracolsep{\fill}} l l l l}
        \toprule
        对象 & 形式 & 关键字段 & 用途 \\
        \midrule
        场景模型 & OBJ/GLTF/GLB/STL & id, name, url, up, bbox & 前端均可预览，当前只有OBJ进入仿真。 \\
        航点 & JSON & x, y, z & Pinia中编辑，生成轨迹前可反复修改。 \\
        轨迹 & 文本 & id, altitude, path & 由航点生成，供配置引用。 \\
        配置 & XML & platform, scanner, speed, angle, frequency & 供HELIOS++读取。 \\
        任务 & 内存对象 & task\_id, status, logs, result & 当前随服务进程生命周期存在。 \\
        结果 & XYZ/LAS/LAZ & X, Y, Z, 可选intensity & 保存原始文件并产生统计和显示抽样。 \\
        \bottomrule
    \end{tabular*}
    \label{tab:dataobjects}
\end{table}

后端用Pydantic进行结构校验，资源不存在返回404，参数或格式错误返回4xx，子进程失败保留日志。前端在上传、加载、请求和连接处捕获异常并给出可操作提示。公网部署时不应暴露服务器绝对路径和内部堆栈。

\section{数据分析}

\subsection{数据对象与组织}

后端请求模型包括TrajectoryRequest、ConfigRequest、SimulationRunRequest和SimulationStatus。前端scene store管理模型和点云选项，waypoints store管理航点，simulation store管理任务、日志、结果与参数。大型Three.js对象由单例useThreeScene管理，store只保存业务镜像，避免响应式系统深度代理图形对象。

\subsection{坐标、统计与性能}

距离单位为米，速度为$\text{m/s}$，扫描频率为$\text{Hz}$，脉冲频率为$\text{kHz}$，扫描角为半角度数。系统当前不读取EPSG，也不进行投影转换，因此模型和轨迹必须预先处于同一以米为单位的坐标系；后续应增加CRS元数据和pyproj转换。

设点数为$n$、高程为$z_i$，则

\begin{equation}
    \bar{z} = \frac{1}{n}\sum_{i=1}^{n} z_i, \qquad
    \sigma_z = \sqrt{\frac{1}{n}\sum_{i=1}^{n} (z_i - \bar{z})^2}.
\end{equation}

中位数、P5和P95描述主体分布并降低极端值影响。若$n > m$且$m = 1,000,000$，显示索引可取

\begin{equation}
    i_k = \left\lfloor \frac{k(n-1)}{m-1} \right\rfloor, \quad k = 0, \dots, m-1.
\end{equation}

该方法简单确定，但并非空间均匀采样；超大点云仍需八叉树、多分辨率或渐进策略\citep{pointCloudRendering,pointCloudTerrain}。

\subsection{数据类型与GIS集成}

FeHALS支持多种三维数据格式的输入与输出，覆盖从模型导入到仿真结果导出的完整数据链路。

在输入方面，系统支持OBJ、GLTF、GLB和STL四种三维模型格式。其中OBJ格式作为HELIOS++原生支持的场景格式，可通过objloader滤镜直接加载到仿真场景中；GLTF/GLB格式利用其高效的二进制传输特性，在前端加载和预览时表现更优；STL格式则提供了与CAD建模软件的数据互通能力。模型上传后，系统自动检测模型的up轴方向（Y-up或Z-up），并在Three.js场景中统一转换为Z-up坐标系，与HELIOS++的坐标系统保持一致。

在输出方面，系统支持XYZ、LAS和LAZ三种点云格式。XYZ格式为纯文本格式，每行记录一个点的三维坐标（X、Y、Z）和可选的强度值，兼容性好但占用空间较大；LAS/LAZ格式为ASPRS标准点云格式，支持存储全波形信息和分类标签，但依赖于HELIOS++构建时是否正确链接了LASlib库。当前默认使用XYZ格式，并在解析时按五十万行分块读取，避免大文件整体载入内存。

在GIS集成方面，系统当前采用以米为单位的世界坐标系，不涉及投影转换和地理配准，因此输入的模型和导出的点云均需预先处于同一投影坐标系中。这一设计简化了系统的实现复杂度，但限制了其直接处理地理参考数据的能力。后续可通过引入pyproj库实现坐标参考系统（CRS）的元数据管理，支持在WGS84经纬度和投影坐标系之间进行动态转换，从而实现与GIS数据源的无缝对接。

\section{主要技术手段}

\subsection{前端与三维显示}

前端采用Vue 3组合式API构建单页应用，Vite提供开发与构建\citep{vite}，Pinia管理共享状态\citep{pinia}，Axios统一请求管理。Vue 3的组合式API（Composition API）是前端架构的核心，它允许将Three.js场景逻辑、航点交互逻辑、仿真控制逻辑分别封装为独立的组合式函数（Composables），每个函数管理自己的响应式状态和生命周期，避免了选项式API中mixins的命名冲突和逻辑分散问题。以三维场景为例，useThreeScene组合式函数封装了场景（Scene）、相机（Camera）、渲染器（Renderer）、轨道控制器（OrbitControls）和光线投射器（Raycaster）的初始化与配置，其他组件通过调用该函数获取场景实例，无需关心底层细节。

在三维渲染方面，Three.js封装了WebGL的底层接口，提供了场景图、材质系统、几何体、加载器和后处理等完整工具链\citep{threejs}。系统使用Scene对象管理所有三维对象，通过Group实现模型、航点和点云的分组管理。模型加载使用OBJLoader、GLTFLoader和STLLoader，每种加载器处理后处理逻辑统一。点云渲染使用BufferGeometry存储顶点位置和颜色，通过PointsMaterial控制点的尺寸、透明度和着色方式，采用单一缓冲区避免逐点对象开销。场景支持固定色、高度色和强度色三种着色模式，色彩映射通过蓝—青—绿—黄—红的渐变色带实现。WebGL作为无插件硬件加速的浏览器标准，适合Web端地形与点云显示\citep{webgl,pointCloudTerrain,pointCloudWeb}。

在交互方面，系统通过Raycaster实现鼠标拾取，与一个不可见的Z=0平面求交得到航点位置。OrbitControls提供场景的旋转、缩放和平移。航点使用红色球体表示，选中时变为黄色，白色折线连接各航点形成航线。拖拽航点时临时禁用相机控制，以防止两者冲突。

点云渲染经过性能优化：渲染点数上限为一百万点，超出时使用等间隔索引抽样；颜色和尺寸参数通过滑块实时调整，触发PointsMaterial的更新而非重新创建几何体；后续可将点云解析放入WebWorker线程，并引入八叉树LOD或渐进渲染以支持千万级点云\citep{pointCloudRendering}。

\subsection{后端与HELIOS++}

FastAPI提供异步路由与OpenAPI文档\citep{fastapi}，Pydantic在业务执行前完成类型校验\citep{pydantic}，asyncio子进程避免长任务阻塞API，WebSocket负责连续日志\citep{WebSocket}。REST有利于资源接口解耦\citep{fielding2000rest}，WebSocket适合持续消息。

后端架构分为四个层次：路由层（routers）定义API端点和HTTP方法映射；模式层（schemas）基于Pydantic定义请求和响应的数据结构；服务层（services）实现业务逻辑，包括模型管理、轨迹生成、配置生成、仿真执行和结果解析；配置层（config）管理文件路径、参数范围和运行时配置。各层之间的依赖关系清晰，便于单元测试和功能扩展。

HELIOS++使用C++实现，基于光线追踪技术模拟激光脉冲在地物表面的散射和返回过程\citep{heliosPlusPlus,heliosPlusPlus2}。其核心算法通过CPU光线追踪并行计算每个激光脉冲与场景几何的交点，根据表面材质属性计算回波能量，生成全波形或离散回波信号。系统支持多类场景、平台和扫描器，并允许在物理真实性和计算成本之间进行权衡配置。本系统采用命令行子进程集成方式，通过Python的\texttt{asyncio.create\_subprocess\_exec}在子进程中启动helios++可执行文件，将标准输出和错误输出通过管道逐行读取，写入任务缓冲并广播到WebSocket。这种方式便于故障隔离，子进程崩溃不会影响API服务进程。

轨迹生成是后端服务的关键模块之一，它将用户编辑的航点列表转换为HELIOS++可直接读取的.trj文件。生成过程包括：计算各航段的方向向量并转换为ARINC 705规范的偏航角；以目标飞行速度（默认10 m/s）和航段距离计算时间步长；在航点转角处使用重复时间戳实现瞬时转向。配置生成模块将用户参数映射为HELIOS++的survey XML和scene XML，包括平台类型、扫描器型号、脉冲频率、扫描频率、扫描角度等参数，并正确设置轨迹文件的列索引映射。

点云结果解析支持LAS/LAZ和XYZ两种格式。LAS/LAZ由laspy读取\citep{laspy}，支持全字段解析；XYZ以文本流逐行读取，每五十万行分块处理，避免大文件内存溢出。NumPy用于统计计算，包括总点数、三维边界、高程均值、标准差、中位数、分位数和40组直方图。Web端LiDAR系统应明确处理、存储和分发边界\citep{lidarWebgis}；开源Web架构可将后端处理与浏览器可视分析组成完整流程\citep{viswebdrone}。

\subsection{协作、测试与改进}

仓库通过GitHub分支、Pull Request和审查协作；README说明安装，CHANGELOG记录演进，TODO区分已实现与计划项。交付前应补充：轨迹、弓字形、参数边界、XML转义、XYZ解析和统计单元测试；所有REST路由的接口测试；小型OBJ场景的HELIOS++集成测试；模型增删、航点拖拽、日志重连、点云着色的前端测试；10万、100万及更大点云的耗时、帧率和内存测试。

后续应优先完成HELIOS++环境诊断、统一错误结构与自动测试；再将任务、配置和资源关系持久化并引入任务队列；随后增加CRS、空间索引、LOD、覆盖热力图和仿真动画；最后完成Docker、反向代理、鉴权、HTTPS、配额与监控。当前FeHALS已打通\enquote{可视化规划—配置—仿真—日志—点云分析}的核心链路，适合作为课程设计和HELIOS++入门工作台。

% 以下是预留的图/表位置，待后续补充具体内容
%
% \begin{figure}[pos=!htbp]
%     \centering
%     \includegraphics[width=0.85\textwidth]{figs/scan_pattern_comparison.pdf}
%     \caption{不同飞行参数下的扫描线密度对比。左：低扫描频率/高速度导致稀疏离散线；右：高扫描频率/低速度实现密集覆盖。}
%     \label{fig:scan_pattern}
% \end{figure}
%
% \begin{figure}[pos=!htbp]
%     \centering
%     \includegraphics[width=0.85\textwidth]{figs/trajectory_yaw_example.pdf}
%     \caption{航迹偏航角计算示意图。相邻航点间的方向向量经atan2(dx, dy)转换为ARINC 705偏航角，转角处使用重复时间戳实现瞬时转向。}
%     \label{fig:yaw_calculation}
% \end{figure}
%
% \begin{table}[pos=!htbp]
%     \centering
%     \caption{不同扫描参数配置下的点云质量对比。}
%     \begin{tabular*}{\textwidth}{@{\extracolsep{\fill}} l l l l}
%         \toprule
%         飞行高度(m) & 扫描角(°) & 脉冲频率(kHz) & 扫描频率(Hz) & 条带宽度(m) & 线间距(m) \\
%         \midrule
%         100 & 30 & 200 & 50 & 53.6 & 0.20 \\
%         100 & 60 & 200 & 50 & 115.5 & 0.20 \\
%         200 & 30 & 200 & 50 & 107.2 & 0.20 \\
%         100 & 30 & 100 & 20 & 53.6 & 0.50 \\
%         \bottomrule
%     \end{tabular*}
%     \label{tab:scan_params}
% \end{table}
%

\FloatBarrier

\newpage
\bibliographystyle{gbt7714-numerical}
\bibliography{cas-refs}

\end{document}